% sections/setting.tex
\section{Setting}
\label{sec:setting}

\subsection{The contest}
\label{sec:contest}

Phase~1 (Round~1) of the AWS Trainium Frontier challenge \citep{trainiumfrontier2026rules} asked teams to train a
language model from scratch with the organiser's kit\footnote{Provided by the organisers to participants. We
do not redistribute the kit, its data loader or its tokenizer, and our recipes import them.} on a
\texttt{trn2.3xlarge} instance \citep{aws_trn2_arch,aws_trainium2_arch}. The rules that shape everything below are
these.
\begin{itemize}
  \item \textbf{A charged-time budget.} Training may use \val{budget-s}~s (30~minutes) of wall-clock time,
    excluding start-up and compilation \citep{trainiumfrontier2026rules}. The kit runs eight data-parallel ranks
    under \texttt{torchrun}.
  \item \textbf{One file.} A submission is a single \texttt{train.py}, run inside the organiser's harness. The
    submission portal rejected files of \val{size-limit} bytes or more; this limit is not stated in the published
    rules. Our final file is \val{k82-bytes} bytes.
  \item \textbf{The metric.} The score is validation bits per byte ``on the pinned validation shard, computed over
    \textasciitilde20M tokens'' \citep{trainiumfrontier2026rules},
    \begin{equation}
      \label{eq:bpb}
      \text{\valbpb{}} \;=\; \frac{1}{\ln 2}\,\frac{\sum_{t} -\ln p_\theta(x_t \mid x_{<t})}{\text{number of UTF-8 bytes}},
    \end{equation}
    which is independent of the tokenizer; lower is better. \Cref{eq:bpb} is also what our rehearsals compute. The kit also ships public validation
    data, which teams may score locally. Uploads are scored only by the organisers.
  \item \textbf{Deadline and standings.} Round~1 closed at \val{rules-deadline}. The rules limit each team to one
    entry \citep{trainiumfrontier2026rules}; \cref{tab:uploads} lists the official score of each of our uploads. Our
    last upload, \upl{K82s7}, was made at \val{k82s7-upload}, inside the deadline. We report our best scored upload,
    \upl{K82s4} (\val{best-official-long}); our last upload, a shuffle-salt draw of the same recipe, differs from it
    by \val{salt7-vs-salt4}. The top ten teams advanced to Round~2.
\end{itemize}

\paragraph{The evaluation text.} The rules do not say whether the pinned shard is disjoint from the public
validation data. The size of our offset (\cref{sec:oracle}) is consistent with the official score being computed on
text like the first \val{rules-shard-tokens} public validation tokens: the same model scores \val{text-gap}\docnum{}
worse on 20M public tokens than on the first 2M. If the shards overlap, our rehearsals scored part of the official
text, and choosing seeds and shuffle salts by their rehearsal score (\cref{sec:salts}) was selection on part of the
evaluation set. We treat that possibility as a threat to validity throughout rather than rule it out.

\paragraph{Terms.} A \emph{micro-batch} is 65{,}536 tokens across the eight ranks, and $k$ is the number of
micro-batches accumulated per optimizer step. A \emph{step} is one optimizer update. A \emph{rehearsal} is a
full-length run of an exact upload file on our own Trainium instance, scored on the first \val{eval-tokens} tokens of
the public validation data. An \emph{upload} is an officially scored submission, named by our recipe identifier
(\upl{K82s4} is recipe K82 with shuffle salt~4). A \emph{chip} is one of our separate \texttt{trn2.3xlarge}
instances. Chips A--D produced the run records the simulator is fitted on; chips E, G and H were added on the final
day and produced only rehearsals, which are not in the released run dataset (\cref{sec:limits}). Chips differ in step
time by up to \val{chip-spread}\docnum{}, and chip~D's older Neuron runtime (2.33.10, against 2.34.10 on chip~C) was
\val{chip-old-runtime}\docnum{} slower, so we compare step counts only within a chip. Symbols are collected in
\cref{tab:notation} (\cref{app:notation}).

\subsection{Determinism}
\label{sec:determinism}

The Neuron compiler turns the training step into static graphs \citep{aws_neuron_docs}. Two different properties
matter, and we have different evidence for each.
\begin{enumerate}
  \item \textbf{Same trajectory, given the same steps.} Given the same file, seed and data order, the optimizer
    trajectory should be identical. Our direct evidence is one pair: a cold rehearsal and a warm rerun of one file on
    19~September scored \val{det-cold}\docnum{} and \val{det-warm}\docnum{}, a difference of \val{det-diff}. That is
    repeatable to $10^{-5}$, not bit-for-bit, and it rests on $n=1$ with older code. The offset statistics do not
    settle the question: a slope of \val{reg-slope} (\cref{sec:oracle}) shows the offset does not depend on the
    recipe, and the within-chip offset SD (\val{off-sd-within}, 95\% CI \val{off-sd-within-lo}--\val{off-sd-within-hi})
    cannot rule out an official run that re-draws the order noise (expected SD about \val{noise-chip-pair},
    \cref{sec:levers}).
  \item \textbf{Same number of steps.} Training stops at a wall-clock target, so the step count depends on the
    host's speed and on run-to-run jitter. Within a seed sweep (one recipe on one chip; the seed changes neither the
    program nor the step time) the step count has an SD of \val{jitter-sd-lo}--\val{jitter-sd-hi} steps and a range
    of \val{jitter-range-lo}--\val{jitter-range-hi} steps (\texttt{runs.jsonl}). Across chips it differs by up to
    \val{chip-spread}\docnum{}.
\end{enumerate}
The oracle of \cref{sec:oracle} needs the first property and is degraded by the second. Determinism is not the
default on GPUs, where atomics in attention backward passes make repeated runs differ
\citep{pytorch_reproducibility,summers2021nondeterminism}; \cref{sec:gpu} measures that for our GPU proxy.

\subsection{The price of a step}
\label{sec:price}

In a fixed time budget, a change that helps per step but costs step time must pay for itself. We price everything
with one constant.

\begin{center}
\fbox{\begin{minipage}{0.94\linewidth}
\textbf{Definition (price of a step).} Near our operating point, at a fixed recipe and batch schedule,
\begin{equation}
  \label{eq:exchange}
  \Delta\text{bpb} \;\approx\; -\kappa\,\ln\frac{S'}{S}, \qquad \kappa=\val{kappa}\ \text{bpb},
\end{equation}
where $S$ is the number of optimizer steps the budget allows. Equivalently, one percent more steps is worth
$\rho=\kappa/100=\val{xr-step-rate}$~bpb. \emph{Effective compute} means optimizer steps at a fixed recipe and
batch schedule; tokens and FLOPs are then proportional to steps. Every ``step equivalent'', compute ratio and gap
in this paper uses $\kappa=\val{kappa}$, with $[\val{kappa}, \val{xr-k-anchor}]$ as the sensitivity range.
\end{minipage}}
\end{center}

\paragraph{Where $\kappa$ comes from.} $\kappa=\val{kappa}$ is the per-step rate we measured by normalising
same-recipe runs between two chips of different speed\docnum{} (the number of comparisons behind it is not recorded
in the repository, so we give it no confidence interval). A second, independent calibration comes from one
\val{anchor-date} run trained for \val{xr-anchor-long}~s instead of \val{xr-anchor-short}~s
(\val{xr-anchor-pct} more steps at a constant 262k-token batch, before the batch warm-up existed), which improved
\valbpb{} by \val{xr-anchor-gain}\docnum{}. One point and one free parameter is a calibration, not a fit: it gives
$\kappa=\val{xr-k-anchor}$. The law is not constant across the campaign; \cref{app:kappa} gives its curvature and
why the surrogate's slope is not independent evidence for it.

\paragraph{Uses and limits.} \Cref{eq:exchange} prices changes of a few percent, which is how we use it in
\cref{sec:findings}; the simulator builds it into its decomposition (\cref{sec:decomp}). Extrapolations to large gaps (\cref{sec:gap}) are indicative only. One consequence is worth
stating early: a change that costs 1\% of step time must improve per-step quality by at least \val{xr-step-rate}~bpb to
break even, so a per-step gain measured at equal steps, on a GPU proxy for instance, is not a win until its step
cost on the chip is known. \Cref{fig:exchange} (\cref{sec:gap}) shows the law with both calibrations.

\subsection{The final recipe}
\label{sec:recipe}

\Cref{tab:recipe} summarises \upl{K82s4}, our best scored file. It descends from the organiser's GPT baseline, and
every setting is a default in the file (\cref{app:recipe} lists the knobs). Its lineage over the last two days was
\upl{K60} (attention-source reuse) $\to$ \upl{K63} (EMA off, cooldown 0.7, lower learning rate in the small-batch
phases) $\to$ \upl{K70a} (row pool) $\to$ \upl{K73s4} (fused optimizer, a deeper execution queue, a lower MLP
output-projection learning rate, salt~4) $\to$ \upl{K77a} (EMA blend 0.6) $\to$ \upl{K82s4} (EMA prewarm). \Cref{sec:steps} breaks down where
the campaign's extra optimizer steps came from.

\begin{table}[tbp]
  \centering
  \caption{\textbf{The final recipe, \upl{K82s4}.} About \val{params}\docnum{} parameters; every value is a default
    of \texttt{recipes/K82s4/train.py} (full knob list in \cref{tab:hyper}).}
  \label{tab:recipe}
  \small
  \begin{tabular}{@{}p{2.2cm}p{13.6cm}@{}}
    \toprule
    Component & Setting \\
    \midrule
    Model & 9 blocks, width 1{,}024; four 256-wide query heads sharing one key/value head (GQA); value embeddings
      fed to blocks 0 and 8; LeakyReLU(0.35)$^2$ MLP ($4\times$); rotary embeddings with a key offset (radius 0.3);
      blocks 6--8 reuse block~5's attention source (\knob{FF_ATTN_SRC=5:6,7,8}); logit soft cap 15. \\
    Objective & Next-token loss plus a 3-token multi-token-prediction loss with weights 1/0.5/0.25, faded to plain
      next-token prediction in eight stages by 45\% of the run. \\
    Optimizers & Muon (5 Newton--Schulz iterations, fused update) for matrices; AdamW for embeddings, scalars and
      the head; MLP output projection at learning rate $\times0.8$. \\
    Schedule & Charged-time target 1{,}795~s; batch warm-up 65k $\to$ 131k $\to$ 262k tokens per step
      (\knob{FF_ACCUM_SCHED=1:0.06,2:0.2,4}) with learning rate $\times0.35$ and $\times0.6$ in the small-batch phases;
      20 warm-up steps; linear cooldown over the last 70\% of the run. \\
    Data & Stream-row loader; from micro-batch \val{rowpool-from} on, rows are drawn without replacement from a pool
      of \val{rowpool-size} micro-batches (shuffle salt 4). \\
    Final weights & 0.6 blend of live weights and an EMA updated every 32 steps; the EMA update is compiled
      during the excluded start-up (\knob{FF_EMA_PREWARM}). \\
    Seed & 73. \\
    \bottomrule
  \end{tabular}
\end{table}
